\section{Crop-label alignment}
\label{sec:kappa}

\subsection{Definition}

A labeller produces, for each crop $c$ and each pixel $p$ inside it, a label
$L_c(p)$. When it reads only the pixels inside the crop, two crops covering the same
source pixel can assign it different labels. A predictor $f$ maps a pixel and the
crop it is read in to a class, written $f(p \mid c)$.

For a source pixel $p$, let $C(p)$ be the set of evaluated crops covering it and
supplying a valid label, and
$m(p) = |C(p)|$. The \emph{artefact set} is
\begin{equation}
\A = \{\, p : \exists\, c, c' \in C(p) \text{ with } L_c(p) \neq L_{c'}(p) \,\}.
\label{eq:aset}
\end{equation}
Off $\A$ the labeller is locally single-valued and there is nothing to detect. A
pixel covered by one crop can never enter $\A$, so $m(p) \geq 2$ throughout.

For $p \in \A$ and $c \in C(p)$, with $\ind{\cdot}$ the indicator, define
\begin{align}
\own(p,c) &= \ind{f(p \mid c) = L_c(p)}, \label{eq:own}\\
\oth(p,c) &= \frac{1}{m(p)-1} \sum_{c' \neq c} \ind{f(p \mid c) = L_{c'}(p)},
\label{eq:oth}
\end{align}
where here and below sums over $c$ or $c'$ run over $C(p)$. \Cref{eq:own} asks
whether the prediction matches the label of the crop being read, \cref{eq:oth} how
often it matches another covering crop's label. Averaging those differences over
\emph{source pixels},
\begin{equation}
\kap = \frac{1}{|\A|} \sum_{p \in \A} \frac{1}{m(p)}
       \sum_{c \in C(p)} \big[\, \own(p,c) - \oth(p,c) \,\big].
\label{eq:kappa}
\end{equation}
Each $p \in \A$ counts once. Averaging over $(p,c)$ instances instead gives
$\kapcr$, the \emph{crop-read} weighting, in which a pixel counts once per covering
crop; on $\A$, coverage runs from 2 to over 100, so the two differ. We report
\cref{eq:kappa} throughout and $\kapcr$ in the supplement. This statistic lies in
$[-1,1]$, is conditional on $p\in\A$, and is unrelated to Cohen's kappa despite the
shared symbol. We therefore report its disagreement prevalence
$P(\A)=|\A|/|\{p:m(p)>0\}|$ over evaluation-valid covered source pixels alongside
the primary benchmark results: $\kap$ describes alignment severity where labels
conflict, not how often conflict occurs.

\subsection{The null is exact}

\begin{proposition}
\label{prop:null}
Let $\A$ be non-empty and let $f$ be crop-invariant on $\A$, meaning
$f(p \mid c) = f(p)$ for every $p \in \A$ and every $c \in C(p)$. Then
$\kap = 0$ exactly.
\end{proposition}

\begin{proof}
Fix $p \in \A$, write $m = m(p)$ and $y = f(p)$, which by hypothesis does not depend
on the crop. For each class $k$ let $n_k = |\{ c \in C(p) : L_c(p) = k \}|$, so that
$\sum_k n_k = m$. Summing \cref{eq:own} over the crops covering $p$,
\begin{equation}
\sum_{c} \own(p,c) = \sum_{c} \ind{y = L_c(p)} = n_y .
\label{eq:sumown}
\end{equation}
Summing \cref{eq:oth} and exchanging the order of summation,
\begin{align}
\sum_{c} \oth(p,c)
  &= \frac{1}{m-1} \sum_{c} \sum_{c' \neq c} \ind{y = L_{c'}(p)} \nonumber \\
  &= \frac{1}{m-1} \sum_{c'} \ind{y = L_{c'}(p)} \cdot (m-1) \nonumber \\
  &= n_y ,
\label{eq:sumoth}
\end{align}
because each $c'$ is counted once for every $c$ other than itself, that is $m-1$
times, and the normaliser cancels it.

So \cref{eq:sumown} and \cref{eq:sumoth} agree at every $p \in \A$, and therefore
\begin{equation}
\sum_{c \in C(p)} \big[ \own(p,c) - \oth(p,c) \big] = n_y - n_y = 0
\quad \text{for every } p \in \A .
\label{eq:innerzero}
\end{equation}
\Cref{eq:kappa} is a weighted sum of that inner quantity, so it is $0$; and because
\cref{eq:innerzero} holds pixel by pixel, \emph{any} weighting over $p$ gives $0$,
including $\kapcr$.
\end{proof}

Nothing in the argument uses the number of classes, the label distribution, the
coverage $m(p)$, or any property of the imagery; it needs only that $\A$ is
non-empty, since $\kap$ averages over the source pixels in it. For fixed evaluated
predictions, zero is an algebraic identity rather than a null distribution that must
be simulated or estimated. This says nothing by itself about population uncertainty
across scenes.

\begin{corollary}
\label{cor:dep}
$\kap \neq 0$ certifies that $f$ is crop-dependent on $\A$.
\end{corollary}

The converse does not hold: $\kap$ sums signed contributions, so alignment and
anti-alignment may cancel, and $\kap = 0$ does not certify crop-invariance.
Attributing a positive $\kap$ specifically to the training-label construction needs
the matched control of \Cref{sec:floor}. So the exact zero is an algebraic identity
and a check on the code, a nonzero $\kap$ is directional, and our primary comparative
conclusions use the treated arm minus its matched control.

\subsection{Why the measured control is not zero}
\label{sec:floor}

A convolutional network reading crops is not a crop-invariant predictor. Zero
padding at the crop border, receptive fields truncated at the window edge, and any
normalisation computed over the crop all make $f(p \mid c)$ depend on $c$ near
boundaries. This is not our observation: convolutional layers are known to read
absolute position from padding~\cite{kayhan2020translation, islam2020position} and to
develop border artefacts because of it~\cite{alsallakh2021mind}. We measure the
consequence directly as $\Omega$, the probability that two crops covering the
same pixel predict differently. Writing $P$ for the pixels with $m(p) \ge 2$,
\begin{equation}
\Omega = \frac{\sum_{p \in P} \sum_{c \neq c' \in C(p)}
        \mathbf{1}[\, f(p \mid c) \neq f(p \mid c') \,]}
       {\sum_{p \in P} m(p)\big(m(p) - 1\big)}.
\label{eq:omega}
\end{equation}
The inner sum is over ordered crop pairs $(c,c')$; the denominator therefore has
$m(p)(m(p)-1)$ terms for pixel $p$.
It sums over every covered pixel rather than over $\A$, since it is a property of
the model and is defined where the labels agree, and weights crop pairs uniformly
rather than source pixels. It is shift consistency~\cite{zhang2019shiftinvariant} on
overlapping windows, and reads $0.031$ for optical-only sea-ice models on repaired labels and
$0.0108$ for flood models on the published ones.

How far this lifts $\kap$ off zero depends on how crop-invariant the labels really
are, and the two benchmarks answer differently. On flood labels regenerated from the
published per-event threshold, which are crop-invariant exactly, a trained network
reads $\kap = -0.0008$, empirically near the theoretical zero. On sea ice the
optical-only repaired control reads $+0.0097$. That is consistent with the labels still
inheriting a per-scene normalisation a $128 \times 128$ window can partly
recover, though we have not isolated that from the other crop effects above.

So the floor is not convolution alone: $\Omega$ is positive
on both benchmarks while $\kap$ reads $-0.0008$ on one and $+0.0097$ on the other,
which a resolution floor could not do. Crop-invariant training labels do not on their
own hold $\kap$ at zero, and \cref{prop:null} does not say they do: it constrains
predictors, not labels. The floor is an interaction of labels, model, crop grid and
imagery, which is why every attribution claim below uses a contrast between a
crop-dependent arm and its matched control on the same grid; raw values are shown to
expose the floor.
